% Riwayat singkat penulis, wajib menurut Permendikbudristek 22/2022 Lampiran III
% Tabel 3.
\chapter{Riwayat Penulis}

Alfa Yohannis adalah dosen Universitas Pradita di Tangerang. Ia menyelesaikan gelar Doctor of Philosophy bidang ilmu komputer di University of York, Magister Teknik bidang Informatika di Institut Teknologi Bandung, dan Sarjana Teknik bidang Teknik Elektro di Universitas Diponegoro. Sebelumnya ia menjadi \textit{research fellow} di University of York, KTP \textit{associate} di Leonardo Ltd, dan dosen di Kalbis Institute.

Bidang keahliannya meliputi rekayasa perangkat lunak berbasis model, sistem informasi, gamifikasi, sistem terdesentralisasi, dan pendidikan ilmu komputer. Ia mengampu mata kuliah Advanced Software Engineering and DevOps yang menjadi dasar buku ini.

Karya terbitannya antara lain \textit{Turning Models Inside Out} (2017), \textit{Towards Model-Based Bias Mitigation in Machine Learning} (2022), dan \textit{LLM-Assisted TOGAF Preliminary Phase for Student Digital Twins} (2025). Penelitian terakhirnya mengukur jejak karbon dan kinerja penyimpanan data citra deret waktu pada PostgreSQL, MongoDB, dan MinIO.
